\documentclass[11pt]{article}

\usepackage[T1]{fontenc}
\usepackage{lmodern}
\usepackage{amsmath,amssymb,amsthm}
\usepackage{booktabs}
\usepackage{microtype}
\usepackage[hidelinks]{hyperref}

\newtheorem{theorem}{Theorem}
\newtheorem{lemma}[theorem]{Lemma}
\newcommand{\Q}{\mathbb{Q}}
\newcommand{\Zf}{\mathcal{Z}^{\mathrm f}_*}

\title{Formal Multiple \(q\)-Zeta Injectivity Through Weight Eight\\
       via an Exact Rank Certificate}
\author{Xiaojian Zeng\\
\small PBC School of Finance, Tsinghua University, Beijing 100190, China\\
\small \texttt{xiaojianzeng@foxmail.com}}
\date{9 September 2026}

\begin{document}
\maketitle

\begin{abstract}
For the Schlesinger--Zudilin model of formal multiple \(q\)-zeta values,
Mozgovoy records a conjecture that quotienting the admissible quasi-shuffle algebra by
the ideal of duality relations gives the algebra of the corresponding formal
\(q\)-series.  We prove the complete weight-at-most-eight instance.  The proof
is computer-assisted but exact: the filtered word space has dimension \(1597\),
an explicit subspace of the duality ideal has rank at least \(1432\), and the
first \(165\) coefficient functionals have rank at least \(165\).  Nonzero
integer minors, certified modulo two primes, establish both lower bounds.
Rank--nullity then identifies the finite coefficient kernel with the displayed
ideal subspace.  The result makes no assertion in weights above eight.
\end{abstract}

\section{Statement and context}

We work over \(\Q\) and regard \(q\) as a formal variable.  Put
\[
 z_i=\frac{q^i}{(1-q)^i}\qquad(i\geq 0),
\]
and use the same symbols for letters satisfying \(z_i=x^iz_0\) in the free
associative algebra \(H=\Q\langle x,z_0\rangle\).  The admissible space
\[
 H^0=\Q 1\oplus xHz_0
\]
has as a basis the empty word and the words
\(z_{s_1}\cdots z_{s_k}\) with \(s_1\geq1\) and
\(s_2,\ldots,s_k\geq0\).  It carries the quasi-shuffle product
\cite{Hoffman2000}
\begin{align*}
 1*u&=u*1=u,\\
 (z_a u)*(z_b v)
 &=z_a(u*z_bv)+z_b(z_au*v)+z_{a+b}(u*v).
\end{align*}
Define the anti-involution \(\tau\) by
\[
 \tau(uv)=\tau(v)\tau(u),\qquad \tau(x)=z_0,
 \qquad \tau(z_0)=x,
\]
and let
\[
 I=(\tau(u)-u:u\in H^0)_*
\]
be the full ideal in the quasi-shuffle algebra, not merely the linear span of
its displayed generators.  We write \(\Zf=H^0_*/I\).

For an admissible word set
\[
 Z(z_{s_1}\cdots z_{s_k})=
 \sum_{n_1>\cdots>n_k>0}
 \prod_{j=1}^k
 \left(\frac{q^{n_j}}{1-q^{n_j}}\right)^{s_j},
 \qquad Z(1)=1.
\]
This is a well-defined element of \(\Q[[q]]\).  The map \(Z\) is a
quasi-shuffle algebra homomorphism, and the duality theorem gives
\(Z(\tau u)=Z(u)\) \cite[Theorem~3.5]{Mozgovoy2026}.  It therefore descends to
a surjection from \(\Zf\) onto its image.

Give a word the weight
\[
 \operatorname{wt}(z_{s_1}\cdots z_{s_k})
   =\sum_{j=1}^k\max(s_j,1),\qquad \operatorname{wt}(1)=0,
\]
and give \(\Zf\) the induced filtration.  Conjecture~3.6 in Mozgovoy asks for
the descended map to be an isomorphism in every weight
\cite{Mozgovoy2026}.  A closely related stuffle/duality quotient and partial
structural results appear in \cite{Brindle2025}.  Takeyama reported that
double-shuffle and resummation-duality relations account for the relations in
his modified \(q\)-MZV computation through weight seven
\cite[p.~2]{Takeyama2013}.  That computation does not itself give the exact
certificate below for the present filtered quotient.

Our result is the following bounded assertion.

\begin{theorem}\label{thm:main}
The canonical map induced by \(Z\) is injective on \(W_8\Zf\).  Equivalently,
for every \(a\in W_8H^0\),
\[
 Z(a)=0\quad\Longrightarrow\quad a\in I.
\]
\end{theorem}

The comparison above records the closest results found in the cited sources;
it is not an exhaustive priority claim.

\section{The two exact matrices}

Let \(V=W_8H^0\).  The generating function for the first, positive letter is
\(t/(1-t)\).  A later letter has generating function
\(2t+t^2/(1-t)\), since both indices \(0\) and \(1\) have weight one.
Consequently the generating function for nonempty admissible words is
\[
 \frac{t/(1-t)}{1-(2t+t^2/(1-t))}
 =\frac{t}{1-3t+t^2}.
\]
The word counts in exact weights \(0,1,\ldots,8\), including the empty word
in weight zero, are
\[
 1,1,3,8,21,55,144,377,987.
\]
Thus
\begin{equation}\label{eq:dim}
 \dim_{\Q}V=1597.
\end{equation}

Fix the ordering by increasing weight and then the recursive index ordering
used in the ancillary certificate.  Define
\begin{equation}\label{eq:J}
 J=\operatorname{span}_{\Q}
 \left\{v*(\tau(u)-u):
 u,v\text{ basis words},\ 
 \operatorname{wt}(u)+\operatorname{wt}(v)\leq8\right\}.
\end{equation}
Since \(\tau\) preserves weight and
\[
 \max(a+b,1)\leq\max(a,1)+\max(b,1)
 \qquad(a,b\geq0),
\]
every vector in \eqref{eq:J} is supported in \(V\).  It also lies in \(I\).
For the actual enumeration, \(z_{s_1}\cdots z_{s_k}\) is encoded by the
binary word \(1^{s_1}0\cdots1^{s_k}0\); reversing the bits and interchanging
zero and one implements \(\tau\) directly.
After zero vectors are removed, these generators give \(4636\) integer rows
in the \(1597\)-word basis.  Let \(R\) be the resulting
\(4636\times1597\) matrix, so that \(J\) is its rational row space.

The second integer matrix records initial series coefficients:
\begin{equation}\label{eq:C}
 C_{n,w}=[q^n]Z(w),\qquad 0\leq n\leq164,
\end{equation}
where \(w\) runs over the same word basis.  Hence \(C\) is
\(165\times1597\).

For completeness, the entries of \(C\) can be constructed without numerical
evaluation.  For \(s\geq0\) and \(m\geq1\), put
\[
 F_{s,m}(q)=
 \begin{cases}
 1,&s=0,\\[2pt]
 \left(q^m/(1-q^m)\right)^s,&s>0.
 \end{cases}
\]
For a possibly nonadmissible suffix \(\boldsymbol{s}=(s,\boldsymbol{t})\),
let \(H_{\boldsymbol{s}}(M)\) be the strict nested sum with largest index at
most \(M\).  Then
\begin{equation}\label{eq:recurrence}
 H_{(s,\boldsymbol{t})}(M)
 =H_{(s,\boldsymbol{t})}(M-1)
  +F_{s,M}H_{\boldsymbol{t}}(M-1),
\end{equation}
with \(H_{\varnothing}(M)=1\) and
\(H_{\boldsymbol{s}}(0)=0\) for nonempty \(\boldsymbol{s}\).  For \(s>0\),
\begin{equation}\label{eq:factorcoeff}
 [q^{m\ell}]F_{s,m}(q)=\binom{\ell-1}{s-1}
 \quad(\ell\geq s),
\end{equation}
and all other coefficients vanish.  Equations
\eqref{eq:recurrence}--\eqref{eq:factorcoeff}, truncated after degree \(164\),
therefore use integer arithmetic only.  Because the first index of an
admissible word is positive, a largest summation index exceeding \(164\)
cannot affect any retained coefficient.  Thus \eqref{eq:C} is exact, not a
floating-point approximation.

\section{Certificate and proof}

The ancillary files record pivot rows and columns for the following minors.

\begin{center}
\begin{tabular}{@{}rrrrr@{}}
\toprule
prime & size in \(R\) & \(\det(R')\bmod p\)
      & size in \(C\) & \(\det(C')\bmod p\)\\
\midrule
1000003 & 1432 & 795164 & 165 & 372525\\
1000033 & 1432 & 221688 & 165 & 994554\\
\bottomrule
\end{tabular}
\end{center}

Both moduli are prime, and every displayed residue is nonzero.  Each selected
minor is the reduction of an integer minor.  Nonvanishing modulo either prime
therefore implies nonvanishing over \(\Q\), and hence
\begin{equation}\label{eq:ranks}
 \dim_{\Q}J=\operatorname{rank}_{\Q}R\geq1432,
 \qquad \operatorname{rank}_{\Q}C\geq165.
\end{equation}
The second prime is a redundant replication, not an additional assumption.

\begin{proof}[Proof of Theorem~\ref{thm:main}]
For every generator used in \eqref{eq:J}, multiplicativity and duality give
\[
 Z\bigl(v*(\tau(u)-u)\bigr)
 =Z(v)\bigl(Z(\tau(u))-Z(u)\bigr)=0.
\]
Thus \(J\subseteq\ker C\).  Combining rank--nullity,
\eqref{eq:dim}, and \eqref{eq:ranks}, we obtain
\[
 1597=\dim V
 \geq \dim J+\operatorname{rank}C
 \geq1432+165=1597.
\]
All inequalities are equalities, so \(\ker C=J\subseteq I\).  If
\(a\in V\) satisfies \(Z(a)=0\) as a formal series, its coefficients through
degree \(164\) vanish.  Hence \(Ca=0\), whence \(a\in J\subseteq I\).
This is precisely injectivity on the induced filtered piece \(W_8\Zf\).
\end{proof}

The argument does not assume a conjectural quotient dimension.  It also does
not need an a priori description of all of \(I\cap V\): the explicit subspace
\(J\subseteq I\) already saturates the complementary coefficient-rank bound.
In particular, possible presentations of elements of \(I\cap V\) involving
higher-weight terms do not create a gap in the proof.

\section{Reproducibility and limitations}

The ancillary program \texttt{weight8\_certificate.py} enumerates the words,
constructs \(R\) and \(C\), and writes the two modular certificates.  Each
certificate contains the selected row and column indices as well as the four
determinant residues.  A separate verifier, which does not import the producer
program, reconstructs the word ordering and ideal rows and recomputes the
selected coefficient columns using out-of-place suffix tables.  It reproduces
the determinant residues and checks that every one of the \(4636\) ideal rows
is annihilated by all \(165\) coefficient rows modulo each prime.

This separation of generated witnesses from a reconstruction verifier is also
used in other exact finite work; for an application using integer predicates
and separate exact implementations, see
\cite{ZengEtAl2026}.  That citation is methodological and no mathematical
claim from that finite-geometry setting is used here.

The theorem is confined to weights at most eight.  It neither proves
Mozgovoy's Conjecture~3.6 in general nor derives the conjectural Poincar\'e
series for \(\Zf\).  Successful verification of the finite files establishes
the stated matrix identities; the mathematical conclusion additionally uses
the human-readable rank--nullity argument above and Mozgovoy's duality theorem.

\paragraph{Computational and AI assistance.}
The certificate construction used exact integer and finite-field arithmetic
implemented in Python with NumPy, Numba, and python-flint.  A separately
written verification program reconstructed the decisive matrices and minors.
OpenAI Codex assisted with literature search, program development, adversarial
audit organization, and manuscript preparation.  Language-model output is
not used as mathematical evidence: the result rests on the displayed
reduction, the serialized exact minors, and the reconstruction verifier.  No
proof assistant was used.

\paragraph{Data and code availability.}
The ancillary package accompanying this manuscript contains the certificate
generator, the two JSON certificates, the reconstruction verifier, and exact
run instructions.

\raggedright
\bibliographystyle{plain}
\bibliography{references}

\end{document}
